% Generated draft: requires mathematical and source review.
\documentclass[11pt]{article}
\usepackage[margin=25mm]{geometry}
\usepackage{fontspec}
\setmainfont{Latin Modern Roman}
\usepackage{amsmath,amssymb,mathtools,unicode-math}
\setmathfont{Latin Modern Math}
\usepackage{enumitem,needspace,xurl,hyperref}
\setlength{\parindent}{0pt}
\setlength{\parskip}{5pt}
\setlength{\emergencystretch}{3em}
\newcommand{\N}{\mathbb{N}}
\newcommand{\Z}{\mathbb{Z}}
\newcommand{\Q}{\mathbb{Q}}
\newcommand{\R}{\mathbb{R}}
\newcommand{\C}{\mathbb{C}}
\newcommand{\F}{\mathbb{F}}
\newcommand{\A}{\mathbb{A}}
\newcommand{\PP}{\mathbb{P}}
\newcommand{\E}{\mathbb{E}}
\newcommand{\eps}{\varepsilon}
\newcommand{\dd}{\,\mathrm{d}}
\begin{document}
\section*{UnsolvedMath ID: 20003328}
\textbf{Source problem code:} catalogue entry
\subsection{Definitions and mathematical statement}
The canonical text is a transition after Conjecture 39: it reports that E. Lindenstrauss proved ``this case'' in [120], except for escape of limiting mass to a cusp in the noncompact case, and introduces an intermediate problem suggested by Lindenstrauss. The inserted words ``OPEN PROBLEMS 17'' are a running header. The precise conjecture, meaning of ``this case'', and following intermediate problem are absent.

According to the supplied aid, the intended background is arithmetic quantum unique ergodicity for normalized Hecke--Maass eigenfunction measures $|\phi_i|^2\,dV$, with $dV$ normalized to total mass one. It reports convergence to $dV$ for compact congruence arithmetic surfaces, and, in the noncompact setting of the cited result, only local limits $c\,dV$ with $0\leq c\leq1$. The missing assertion is $c=1$. The aid attributes a later nonescape result for Hecke--Maass forms on the modular surface to Soundararajan and says that his argument was stated to apply to congruence subgroups. These scope and literature claims are supplied context, not a newly verified status determination.

The aid's companion criterion can be stated independently. Let $X$ be a locally compact metric space, let $H:X\to[1,\infty)$ be continuous and proper, let $\mu_i,\nu$ be Borel probability measures, and assume
\[
 \mu_i\longrightarrow c\nu\text{ vaguely},\quad 0\leq c\leq1,
 \qquad \sup_i\int_X H^p\,d\mu_i\leq M<\infty
 \quad(p>0).
\]
Then the measures are tight, $c=1$, $\mu_i$ converge weakly to $\nu$ as probabilities, and $\int H^p\,d\nu\leq M$. To formulate the reported rate, let $0\leq q<p$, $|f|\leq L H^q$, and choose continuous cutoffs $\chi_R$ equal to one on $\{H\leq R\}$, zero on $\{H\geq2R\}$, and between zero and one. If, for $R\geq1$,
\[
 \left|\int\chi_R f\,d\mu_i-\int\chi_R f\,d\nu\right|
 \leq D\varepsilon_i R^a,\qquad
 a\geq0,\quad \varepsilon_i\to0^+,
\]
then
\[
 \left|\int f\,d\mu_i-\int f\,d\nu\right|
 \leq D\varepsilon_i R^a+2LMR^{-(p-q)}
 =O\!\left(\varepsilon_i^{(p-q)/(a+p-q)}\right)
\]
after choosing $R$ of order $\varepsilon_i^{-1/(a+p-q)}$. The compact-test estimate and uniform moment are assumptions, not estimates proved here for arithmetic eigenfunctions.

For the cusp normalization mentioned by the aid, a cusp with coordinate $0\leq x<w_j$, $y\geq Y_j$, and hyperbolic area $dx\,dy/y^2$ on a surface of total area $A$ has normalized tail mass $w_j/(AR)$ above height $R\geq Y_j$. Its height-$p$ integral is
\[
 \frac{w_j}{A}\int_{Y_j}^{\infty}y^{p-2}\,dy
 =\frac{w_jY_j^{p-1}}{A(1-p)}\quad(0<p<1),
\]
and diverges for $p\geq1$. Summing these expressions over disjoint cusps gives the corresponding cusp contributions; compact-core contributions depend on the chosen height extension. Thus ordinary cusp height only permits finite target moments with $p<1$. These coordinate conventions clarify the reported companion discussion and do not reconstruct the missing intermediate problem.
\subsection{Short English statement}
This fragment identifies escape of mass as the remaining issue in a cited arithmetic quantum-limit result. A supplied companion criterion says that a uniform positive moment of a proper height prevents escape and can turn compact-test estimates into full convergence rates. Neither the surrounding conjecture nor Lindenstrauss's following intermediate problem is actually stated in the record.
\subsection{Sources}
\begin{itemize}
\item Local UnsolvedMath registry, permanent ID 20003328, exact source code catalogue entry; \texttt{aim-other-notes.json}, notes[134], transition after Conjecture 39.
\item American Institute of Mathematics, \emph{Emerging applications of measure rigidity}, \url{https://aimath.org/WWN/measrigid/measrigid.pdf}. The source attributes the cited case and the suggested intermediate problem to E. Lindenstrauss and uses bibliography label [120], whose full entry is not supplied. Original remarks and literature fields are empty.
\item Supplied research attempt 1: context\_only, problem status not\_a\_problem, rigor label proved, novelty confidence low. It reports the Lindenstrauss/Soundararajan arithmetic QUE qualifications and the height-moment reduction. Machine-generated, not peer reviewed, requiring independent verification before citation. No complete Soundararajan reference or exact eigenfunction-family hypotheses are supplied.
\end{itemize}
\end{document}
